% generated by roam.py handbook export -- не редактировать, править заметки
% 20260918211205-дифференциальныи_критерии_выпуклости_функции_однои_переменнои.org (body abe91812a6bc)

\begin{theorem}[{Дифференциальный критерий выпуклости функции одной переменной}]\label[theorem]{note:2F475EE4-2630-45FD-8DD2-8EA9D48059E4}
  Пусть $I \subset \R$~--- промежуток, $g: I \rightarrow \R$~--- дифференцируемая функция.
  Тогда $g$~--- \hyperref[note:1CA14F2C-637F-41DE-9E97-72075D19E97F]{выпуклая} $\Leftrightarrow$ $\forall s,t \in I\ g(t) \ge g(s) + g^{\prime}(s)(t-s)$
\end{theorem}

\begin{remark}
  Геометрически: график выпуклой функции лежит не ниже любой своей касательной.
\end{remark}

\begin{proof}
  \proofpart{$\Rightarrow$}
  Пусть $g$~--- выпуклая, $s \ne t$ (при $s = t$ неравенство очевидно). Возьмем $\lambda \in (0;1]$, по выпуклости:
  \[
  g(s + \lambda(t-s)) = g((1-\lambda)s + \lambda t) \le (1-\lambda) g(s) + \lambda g(t)
  \]
  Перенесем $g(s)$ влево и поделим на $\lambda > 0$:
  \[
  \frac{g(s + \lambda(t-s)) - g(s)}{\lambda} \le g(t) - g(s)
  \]
  Домножим и поделим левую часть на $t-s$ и устремим $\lambda \rightarrow +0$:
  \[
  \frac{g(s + \lambda(t-s)) - g(s)}{\lambda (t-s)} (t-s) \rightarrow g^{\prime}(s)(t-s)
  \]
  Нестрогое неравенство сохраняется в пределе, значит $g^{\prime}(s)(t-s) \le g(t) - g(s)$.

  \proofpart{$\Leftarrow$}
  Возьмем $s,t \in I$, $\lambda \in [0;1]$ и точку $z = \lambda s + (1-\lambda) t \in I$. Запишем неравенство из условия в точке $z$ дважды:
  \[
  g(s) \ge g(z) + g^{\prime}(z)(s-z)
  \]
  \[
  g(t) \ge g(z) + g^{\prime}(z)(t-z)
  \]
  Умножим первое на $\lambda$, второе на $1-\lambda$ и сложим:
  \[
  \lambda g(s) + (1-\lambda) g(t) \ge g(z) + g^{\prime}(z)(\lambda s + (1-\lambda) t - z) = g(z)
  \]
  Это и есть определение выпуклости: $g(\lambda s + (1-\lambda)t) \le \lambda g(s) + (1-\lambda) g(t)$.
\end{proof}
